\documentclass[12pt,a4paper]{article}
% Drake_preamble.tex - LaTeX Preamble Template
% 作者: 謝慕揚, MD, PhD, FESC
% 
% 使用說明:
% 1. 表格請使用 longtable，不要有垂直分隔線 |
% 2. 表格內換行使用 \newline，不要使用 \\
% 3. 藥物名稱、檢驗、檢查請維持英文
% 4. Reference 格式請使用 \href 加上驗證過的 DOI

% Including essential packages for Chinese typesetting
\usepackage{xeCJK}
\usepackage{fontspec}
% Configuring Chinese font with Noto Serif CJK TC, fallback to SimSun
\IfFontExistsTF{Noto Serif CJK TC}{
  \setCJKmainfont{Noto Serif CJK TC}
  \setCJKsansfont{Noto Serif CJK TC}
  \setCJKmonofont{Noto Serif CJK TC}
}{\setCJKmainfont{SimSun}
  \setCJKsansfont{SimSun}
  \setCJKmonofont{SimSun}
}

% Including other necessary packages
\usepackage{geometry}
\usepackage{titlesec}
\usepackage{enumitem}
\usepackage{xcolor}
\usepackage{colortbl}
\usepackage{tcolorbox}
\usepackage{booktabs}
\usepackage{longtable}
\usepackage{array}
\usepackage{graphicx}
\usepackage{tikz}
\usepackage{fancyhdr}
\usepackage{multicol}
\usepackage{datetime2}
\usepackage{hyperref}
\usepackage{mdframed}
\usepackage{amsmath}
\usepackage{amssymb}
\usepackage{multirow}

\hypersetup{
    colorlinks=true,
    linkcolor=emergency_blue,
    citecolor=emergency_blue,
    urlcolor=emergency_blue,
    pdfborder={0 0 0}
}

% Defining page geometry
\geometry{
  top=2.5cm,
  bottom=2.5cm,
  left=2cm,
  right=2cm,
  headheight=25pt
}

% Adjusting topmargin to compensate for increased headheight
\addtolength{\topmargin}{-10pt}

% Defining colors
\definecolor{emergency_red}{RGB}{186,24,27}
\definecolor{emergency_blue}{RGB}{0,114,188}
\definecolor{emergency_gray}{RGB}{120,120,120}
\definecolor{highlight_yellow}{RGB}{255,250,205}

% Setting title formats
\titleformat{\section}[block]{\Large\bfseries\color{emergency_red}}{\thesection}{10pt}{}
\titleformat{\subsection}[block]{\large\bfseries\color{emergency_blue}}{\thesubsection}{8pt}{}
\titleformat{\subsubsection}[block]{\normalsize\bfseries\color{emergency_gray}}{\thesubsubsection}{6pt}{}

% Defining custom environment for important notes
\newmdenv[
    linecolor=emergency_red,
    backgroundcolor=highlight_yellow,
    linewidth=2pt,
    topline=true,
    bottomline=true,
    leftline=true,
    rightline=true,
    innertopmargin=10pt,
    innerbottommargin=10pt,
    innerleftmargin=10pt,
    innerrightmargin=10pt
]{importantbox}

% Defining note command with tcolorbox
\newcommand{\note}[1]{
    \par
    \noindent
    \begin{tcolorbox}[
        colback=emergency_blue!5!white,
        colframe=emergency_blue,
        title=\textbf{重點提醒},
        fonttitle=\bfseries,
        boxrule=0.5pt,
        arc=3pt,
        left=6pt,
        right=6pt,
        top=6pt,
        bottom=6pt
    ]
    #1
    \end{tcolorbox}
}

% Key findings box
\newcommand{\keyfinding}[1]{
    \par
    \noindent
    \begin{tcolorbox}[
        colback=yellow!10!white,
        colframe=orange!75!black,
        title=\textbf{核心發現摘要},
        fonttitle=\bfseries,
        boxrule=0.5pt,
        arc=3pt
    ]
    #1
    \end{tcolorbox}
}

% Defining custom date format for ISO 8601
\DTMnewdatestyle{iso}{
  \renewcommand{\DTMdisplaydate}[4]{\number##1-\number##2-\number##3}
  \renewcommand{\DTMDisplaydate}[4]{\number##1-\number##2-\number##3}
}
\DTMsetdatestyle{iso}
\newcommand{\mydate}{\DTMtoday}


% Custom header for this document
\pagestyle{fancy}
\fancyhf{}
\fancyhead[L]{\textcolor{emergency_red}{Intensive Care Medicine 教學講義}}
\fancyhead[R]{\textcolor{emergency_blue}{Metabolic Acidosis Buffering}}
\fancyfoot[C]{\textcolor{emergency_gray}{\thepage}}
\renewcommand{\headrulewidth}{1pt}
\renewcommand{\headrule}{\hbox to\headwidth{\color{emergency_red}\leaders\hrule height \headrulewidth\hfill}}

\begin{document}

\begin{center}
{\LARGE\bfseries\textcolor{emergency_red}{Intensive Care Medicine}}\\[0.3cm]
{\Large\textbf{Understanding the Disease}}\\[0.5cm]
{\Large\textbf{Understanding Buffering of Metabolic Acidosis in Critical Illness}}\\[0.3cm]
{\large 了解重症病人代謝性酸中毒的緩衝作用}\\[0.5cm]
{\normalsize Intensive Care Med 2024;50:1925-1928. DOI: 10.1007/s00134-024-07656-5}\\[0.3cm]
{\normalsize 整理：謝慕揚 MD, PhD, FESC \quad 日期：\mydate}
\end{center}

\keyfinding{
\textbf{核心訊息：}相同程度的 metabolic acidosis（$\Delta$SBE）在不同病人可能導致兩倍或一半的 acidaemia（$\Delta$pH），取決於 extracellular fluid 的 buffer power（$\beta$）。

\textbf{實用重點：}
\begin{itemize}
    \item \textbf{Buffer power 公式：}$\beta$ = 2.3$\cdot$[HCO$_3^-$] + 1.43$\cdot$[Hb]/3 + 0.17$\cdot$[Alb]
    \item \textbf{正常值：}$\beta$ $\approx$ 70 mMol/L（約 80\% 來自 bicarbonate）
    \item \textbf{臨床意涵：}低 $\beta$ 病人應考慮 balanced solutions、RBC transfusion 或 sodium bicarbonate
\end{itemize}
}

\tableofcontents
\newpage

%==============================================================
\section{文獻概述}
%==============================================================

本篇為 Intensive Care Medicine 的「Understanding the Disease」系列文章，由英國 Guy's and St Thomas' Hospital 的 Giosa 等人撰寫，解釋 buffer power 的概念及其在重症病人的臨床相關性。

\subsection{作者資訊}
\textbf{通訊作者：}Lorenzo Giosa, MD\\
\textbf{機構：}Department of Critical Care Medicine, Guy's and St Thomas' NHS Foundation Trust, St Thomas' Hospital, London, UK

%==============================================================
\section{背景知識}
%==============================================================

\subsection{Metabolic Acidosis 定義}

\begin{itemize}
    \item 定義為 extracellular fluid (ECF) 的 base excess 下降，或 standard base excess (SBE) 下降
    \item 在重症病人常見，可由以下原因引起：
    \begin{itemize}
        \item Sepsis
        \item Acute kidney injury (AKI)
        \item Normal saline 輸液
    \end{itemize}
\end{itemize}

\subsection{代償機制}

\begin{itemize}
    \item \textbf{呼吸代償：}需數小時才能完全建立
    \item \textbf{腎臟代償：}需數天才能完全建立
    \item \textbf{急性期：}pH 變化主要取決於身體緩衝酸負荷的內在能力
\end{itemize}

%==============================================================
\section{Buffer Power（緩衝能力）}
%==============================================================

\subsection{定義}

\note{
\textbf{Buffer Power ($\beta$) 定義：}
$$\beta = \frac{\Delta SBE}{\Delta pH}$$

意義：需要多少 non-volatile acid 才能使 pH 改變 1 個單位
}

\subsection{緩衝系統組成}

急性 metabolic acidosis 的緩衝主要發生在 extracellular fluid（plasma、red blood cells、interstitial fluid）。

\begin{longtable}{p{4cm}p{10cm}}
\toprule
\textbf{緩衝類型} & \textbf{組成成分} \\
\midrule
\endfirsthead

\toprule
\textbf{緩衝類型} & \textbf{組成成分} \\
\midrule
\endhead

\bottomrule
\endfoot

\textbf{Carbonic Buffer} & Bicarbonate (HCO$_3^-$) \\
\midrule
\textbf{Non-Carbonic Buffers} &
• Hemoglobin (Hb)\newline
• Albumin (Alb)\newline
• Inorganic phosphate (Pi)\newline
• 2,3-bisphosphoglycerate (2,3-DPG) \\
\bottomrule
\end{longtable}

\subsection{Buffer Power 計算公式}

\note{
\textbf{ECF Buffer Power 公式：}
$$\beta_{mMol/L} = 2.3 \cdot [HCO_3^-]_{mMol/L} + 1.43 \cdot \frac{[Hb]_{g/dL}}{3} + 0.17 \cdot [Alb]_{g/L}$$

\textbf{正常值（HCO$_3^-$ = 24 mMol/L, Hb = 15 g/dL, Alb = 45 g/L）：}
\begin{itemize}
    \item $\beta$ $\approx$ 70 mMol/L
    \item 約 80\% 來自 bicarbonate
\end{itemize}

\textbf{臨床意義：}需要約 70 mMol non-volatile acid 加入 1 L ECF 才能使 pH 下降 1 個單位（在 PCO$_2$ 不變的情況下）
}

%==============================================================
\section{臨床案例模擬}
%==============================================================

三位不同背景的病人發展相同的急性狀況（pneumonia 後發生 metabolic acidosis）：

\subsection{基線 Buffer Power 比較}

\begin{longtable}{p{2.5cm}p{2cm}p{2cm}p{2cm}p{2.5cm}p{2.5cm}}
\toprule
\textbf{病人} & \textbf{HCO$_3^-$} & \textbf{Hb} & \textbf{Alb} & \textbf{基線 $\beta$} & \textbf{特點} \\
\midrule
\endfirsthead

\toprule
\textbf{病人} & \textbf{HCO$_3^-$} & \textbf{Hb} & \textbf{Alb} & \textbf{基線 $\beta$} & \textbf{特點} \\
\midrule
\endhead

\bottomrule
\endfoot

\textbf{COPD} & 32.4 & 17 g/dL & 45 g/L & $\sim$90 & 高 $\beta$ \\
\midrule
\textbf{CKD} & 18.9 & 10 g/dL & 30 g/L & $\sim$53 & 低 $\beta$ \\
\midrule
\textbf{AML} & 24.3 & 6 g/dL & 20 g/L & $\sim$62 & 中等 $\beta$ \\
\bottomrule
\end{longtable}

\subsection{Pneumonia 對 Buffer Power 的影響}

\begin{longtable}{p{3cm}p{4cm}p{3.5cm}p{3.5cm}}
\toprule
\textbf{病人} & \textbf{對 Pneumonia 的反應} & \textbf{HCO$_3^-$ 變化} & \textbf{$\beta$ 變化} \\
\midrule
\endfirsthead

\toprule
\textbf{病人} & \textbf{對 Pneumonia 的反應} & \textbf{HCO$_3^-$ 變化} & \textbf{$\beta$ 變化} \\
\midrule
\endhead

\bottomrule
\endfoot

\textbf{COPD} & Hypercapnia 惡化\newline Respiratory acidosis & 上升至 34.4 & 上升至 $\sim$95 \\
\midrule
\textbf{CKD} & Hyperventilation\newline Respiratory alkalosis & 下降至 18.2 & 下降至 $\sim$51 \\
\midrule
\textbf{AML} & Hyperventilation\newline Respiratory alkalosis & 下降至 23.8 & 下降至 $\sim$60 \\
\bottomrule
\end{longtable}

\subsection{相同 Metabolic Acidosis 的不同結果}

\note{
\textbf{關鍵發現：}所有病人發展相同程度的 acute metabolic acidosis（$\Delta$SBE $\approx$ -12 mMol/L），但由於 buffer power 不同，pH 下降幅度差異顯著：

\begin{longtable}{p{3cm}p{3cm}p{3cm}p{4.5cm}}
\toprule
\textbf{病人} & \textbf{$\beta$} & \textbf{$\Delta$pH} & \textbf{最終 pH} \\
\midrule
\endfirsthead

\bottomrule
\endfoot

\textbf{COPD} & $\sim$95 & -0.15 & 7.10 \\
\textbf{CKD} & $\sim$51 & \textbf{-0.30} & 7.10 \\
\textbf{AML} & $\sim$60 & -0.25 & 7.20 \\
\bottomrule
\end{longtable}

\textbf{結論：}CKD 病人的 pH 下降是 COPD 病人的兩倍（0.30 vs 0.15）！
}

%==============================================================
\section{臨床意涵}
%==============================================================

\subsection{理解 Buffer Power 的重要性}

\begin{itemize}
    \item 有助於預測 metabolic acidosis 病人可能需要的支持程度
    \item 有助於個人化管理策略
\end{itemize}

\subsection{低 Buffer Power 病人的處置建議}

\begin{longtable}{p{4.5cm}p{9.5cm}}
\toprule
\textbf{介入措施} & \textbf{說明} \\
\midrule
\endfirsthead

\toprule
\textbf{介入措施} & \textbf{說明} \\
\midrule
\endhead

\bottomrule
\endfoot

\textbf{Fluid Choice} & 優先選擇 balanced solutions（而非 normal saline）\\
\midrule
\textbf{RBC Transfusion} & 增加 hemoglobin 可增加 non-carbonic buffer power \\
\midrule
\textbf{Sodium Bicarbonate} & 直接增加 carbonic buffer capacity \\
\bottomrule
\end{longtable}

\subsection{Albumin Infusion 的效果}

\begin{itemize}
    \item 增加 non-carbonic buffer power
    \item 但同時因增加 plasma 的 fixed negative charges 而減少 bicarbonate
    \item \textbf{淨效果：}對 $\beta$ 的影響可能可忽略
\end{itemize}

\subsection{Permissive Hypercapnia 的潛在價值}

\begin{itemize}
    \item 延長的 permissive hypercapnia 可增加 HCO$_3^-$
    \item 在 ICU 中，metabolic acidosis 風險高且 non-carbonic $\beta$ 常因 anemia 及 hypoalbuminemia 而下降
    \item Hypercapnia 可能提供額外的緩衝保護
\end{itemize}

%==============================================================
\section{高風險病人特徵}
%==============================================================

\note{
\textbf{低 Buffer Power 高風險病人：}
\begin{itemize}
    \item 低 HCO$_3^-$（如 CKD、chronic metabolic acidosis）
    \item 低 Hemoglobin（如 anemia、血液系統疾病）
    \item 低 Albumin（如 營養不良、肝病、nephrotic syndrome）
\end{itemize}

這些病人對 acute metabolic acidosis 的耐受性較差，相同程度的酸負荷會導致更嚴重的 acidaemia。
}

%==============================================================
\section{Take-Home Message}
%==============================================================

\begin{enumerate}
    \item \textcolor{red}{\textbf{Buffer Power 變異大：}}重症病人的 metabolic acidosis 緩衝能力變異極大
    \item \textcolor{red}{\textbf{相同 $\Delta$SBE，不同 $\Delta$pH：}}相同程度的 acidosis 可導致兩倍或一半的 acidaemia，取決於 ECF 的 buffer power
    \item \textcolor{red}{\textbf{關鍵決定因子：}}HCO$_3^-$、Hemoglobin、Albumin 是主要決定 buffer power 的因素
    \item \textcolor{red}{\textbf{個人化管理：}}了解病人的 buffer power 可幫助預測疾病軌跡並個人化治療
\end{enumerate}

%==============================================================
\section{縮寫對照表}
%==============================================================

\begin{longtable}{p{3cm}p{11cm}}
\toprule
\textbf{縮寫} & \textbf{全名} \\
\midrule
\endfirsthead

\toprule
\textbf{縮寫} & \textbf{全名} \\
\midrule
\endhead

\bottomrule
\endfoot

$\beta$ & Buffer Power (緩衝能力) \\
SBE & Standard Base Excess (標準鹼剩餘) \\
ECF & Extracellular Fluid (細胞外液) \\
Hb & Hemoglobin (血紅素) \\
Alb & Albumin (白蛋白) \\
Pi & Inorganic Phosphate (無機磷酸鹽) \\
2,3-DPG & 2,3-Bisphosphoglycerate \\
COPD & Chronic Obstructive Pulmonary Disease \\
CKD & Chronic Kidney Disease \\
AML & Acute Myeloid Leukemia \\
AKI & Acute Kidney Injury \\
RBC & Red Blood Cell \\
PCO$_2$ & Partial Pressure of Carbon Dioxide \\
\bottomrule
\end{longtable}

%==============================================================
\section{參考文獻}
%==============================================================

\begin{enumerate}
    \item Giosa L, Camporota L, Langer T. Understanding buffering of metabolic acidosis in critical illness. \href{https://doi.org/10.1007/s00134-024-07656-5}{\textit{Intensive Care Med}. 2024;50:1925-1928.}

    \item Berend K. Diagnostic use of base excess in acid-base disorders. \href{https://doi.org/10.1056/NEJMra1711860}{\textit{N Engl J Med}. 2018;378:1419-1428.}

    \item Kraut JA, Madias NE. Metabolic acidosis: pathophysiology, diagnosis and management. \href{https://doi.org/10.1038/nrneph.2010.33}{\textit{Nat Rev Nephrol}. 2010;6:274-285.}

    \item Siggaard-Andersen O. The acid-base status of the blood. Williams \& Wilkins Company, Baltimore; 1974.

    \item Langer T, Brusatori S, Carlesso E, et al. Low noncarbonic buffer power amplifies acute respiratory acid-base disorders in patients with sepsis: an in vitro study. \href{https://doi.org/10.1152/japplphysiol.00787.2020}{\textit{J Appl Physiol}. 2021;131:464-473.}
\end{enumerate}

\end{document}
